\section{Заключение}

В ходе выполнения лабораторной работы с помощью критерия Пирсона была проверена гипотеза о том, что случайная выборка объёма $n=100$ имеет экспоненциальное распределение с параметром $\lambda=1.41$. На основе равномерной сетки из 10 интервалов и уровня значимости $\epsilon=0.05$ были определены:
\begin{itemize}
    \item значение функции отклонения \textit{(статистика Пирсона)}: $\rho(\vec{X}) = 17.8$;
    \item критическое значение $\chi^{2}$-распределения с 9 степенями свободы: $C = 16.92$.
\end{itemize}

В результате проверки, так как $\rho(\vec{X}) \ge C$, была принята альтернативная гипотеза $H_{2}$:
\[
    \mathcal{F} \ne \textrm{Exp}(\lambda)
\]

Для автоматизации вычислений \textit{(расчёта квантилей, частотного распределения и статистики)} использовались скрипты, написанные на языке \textit{Python}. Их листинги приведены в теле отчёта.